\section{Introduction}
\label{sec:intro}

Large language model agents are increasingly used to run and combine scientific software, but the software itself is rarely written with agents in mind: a typical research repository assumes a human will read the README, create an environment by hand, and call functions from a Python shell or a bespoke CLI. Turning such a repository into a tool an agent can call reliably — with a stable interface, an isolated environment, and some evidence that it actually works — is a repetitive engineering task that today falls on a person for every repository.

\textbf{Problem and importance.} As the number of scientific repositories that agents are asked to use grows, hand-wrapping each one does not scale. Recent benchmarks confirm that setting up and executing arbitrary research code is itself a hard, largely unsolved agent capability \citep{super2024,corebench2024,astabench2025,researchenvbench2026} — exactly the step Alembic automates. A system that performs this repo-to-tool conversion automatically directly increases the pool of software an agent-based scientific assistant (e.g., a co-scientist system) can draw on, without additional engineering effort per tool.

\textbf{What Alembic does.} Given nothing but a repository URL, Alembic runs a four-stage agent pipeline inside an ephemeral Docker container — \emph{Explorer} (understand the repo), \emph{Environment} (build an installable venv), \emph{Coder} (write a FastMCP server and tests), and \emph{Validator} (prove the server works, calling a \emph{Debugger} sub-agent on failure) — and on success commits the container to a runnable image that serves the generated tools over MCP. Section~\ref{sec:system} describes the pipeline; Section~\ref{sec:demo} describes what a user sees; Section~\ref{sec:related} positions Alembic against the closest prior systems.

\textbf{Novelty.} Unlike systems that require a human-specified target function and invocation example (ToolMaker; \citealp{toolmaker2025}) or that stop at a smoke-tested MCP scaffold on the local filesystem (ToolRosella/code2mcp; \citealp{toolrosella2026}), Alembic (a) discovers the tool surface itself from the repository rather than from a human-authored task spec, and (b) its final artifact is a self-contained, versioned Docker image serving a live MCP endpoint, validated by actually invoking every generated tool rather than only import-checking the generated code. Section~\ref{sec:related} details these differences axis by axis.
